% E. Decisiones técnicas, incluyendo normas y estándares o justificando su ausencia
% (Guía apdo. 3.E). Extensión orientativa: 6-8 páginas. Es el capítulo técnico central.
% Fuentes de verdad: docs/ARCHITECTURE.md, docs/DATA_SCHEMA.md, README.md (pinout)
\chapter{Diseño del sistema}
\label{cap:diseno}

\section{Arquitectura general}

El sistema se estructura en dos niveles. El nodo de borde, basado en un microcontrolador
ESP32, realiza la adquisición de las señales físicas y su empaquetado. El concentrador
local, implementado sobre un ordenador de placa reducida Raspberry~Pi, aloja el
intermediario de mensajería y el registrador que consolida el conjunto de datos destinado al
entrenamiento del modelo.

La figura~\ref{fig:arquitectura} recoge el flujo completo de la información. Conviene
subrayar la dirección de las flechas: el nodo transmite características ya calculadas y no
señal en bruto, y el concentrador devuelve al nodo un modelo entrenado. El sistema no es por
tanto una cadena lineal de sensor a servidor, sino un ciclo en el que el concentrador
interviene en la fase de entrenamiento y queda prescindible en la de explotación.

\begin{figure}[ht!]
\centering
\begin{tikzpicture}[
  font=\small,
  bloque/.style={draw, rounded corners=2pt, align=center, inner sep=6pt, minimum height=1cm},
  nodo/.style={bloque, fill=black!4, text width=4.4cm},
  hub/.style={bloque, fill=black!8, text width=4.4cm},
  activo/.style={bloque, fill=black!12, text width=2.6cm},
  flecha/.style={-{Latex[length=2mm]}, thick},
  etiq/.style={font=\scriptsize, align=center, inner sep=2pt}
]
  \node[activo] (activo) {Compresor\\(activo de prueba)};
  \node[nodo, right=2.4cm of activo] (nodo)
      {\textbf{Nodo de borde (ESP32)}\\[2pt]
       Adquisición multisensor\\
       Extracción de características\\
       Inferencia embarcada};
  \node[hub, right=2.4cm of nodo] (hub)
      {\textbf{Concentrador (Raspberry Pi)}\\[2pt]
       Intermediario MQTT\\
       Registrador a ficheros CSV\\
       Entrenamiento del modelo};

  \draw[flecha] (activo) -- node[etiq, above] {vibración\\temperatura\\sonido} (nodo);
  \draw[flecha] (nodo) -- node[etiq, above] {MQTT sobre\\Wi-Fi 2,4 GHz} (hub);
  \draw[flecha] (hub.south) -- ++(0,-0.9)
      node[etiq, below] {modelo entrenado (fase de despliegue)} -| (nodo.south);
  \draw[flecha] (nodo.north) -- ++(0,0.8)
      node[etiq, above] {veredicto de estado};
\end{tikzpicture}
\caption{Arquitectura del sistema. El nodo transmite características extraídas, no señal en
bruto; el concentrador entrena el modelo y lo devuelve al nodo, que ejecuta la inferencia.}
\label{fig:arquitectura}
\end{figure}

\section{Adquisición: sensores, plataforma y frecuencias de muestreo}

El cuadro~\ref{table:sensores} recoge los sensores seleccionados y la magnitud que aporta
cada uno al diagnóstico.

\begin{table}[ht!]
\caption{Sensores integrados en el nodo de borde.}
\label{table:sensores}
\centering
\begin{tabular}{|p{3cm}|p{2.6cm}|p{7cm}|}
\hline
\textbf{Sensor} & \textbf{Interfaz} & \textbf{Magnitud y aportación al diagnóstico} \\
\hline
MPU-6050~\cite{mpu6050} & I\textsuperscript{2}C & Aceleración en tres ejes y velocidad angular en tres
ejes; detecta desequilibrios mecánicos y aflojamiento del anclaje. Aporta además una
temperatura interna utilizada como indicador relativo. \\
\hline
DS18B20~\cite{ds18b20} & 1-Wire & Temperatura del entorno o de la tubería, con encapsulado metálico apto
para contacto; permite evaluar el rendimiento térmico del ciclo. \\
\hline
INMP441~\cite{inmp441} & I2S & Nivel acústico digital omnidireccional; la salida digital evita la
inmunidad reducida a la interferencia electromagnética propia de los micrófonos analógicos
en presencia de un motor. \\
\hline
\end{tabular}
\end{table}

La plataforma de procesado quedó condicionada por los buses que exigen esos sensores. Se
seleccionó el ESP32 por integrar en un mismo encapsulado los tres requeridos
(I\textsuperscript{2}C, 1-Wire por software e I2S nativo), conectividad inalámbrica y capacidad
de cómputo suficiente para ejecutar inferencia embarcada~\cite{esp32datasheet}, con un consumo y
un coste compatibles con el despliegue en activos de valor moderado. El
cuadro~\ref{table:requisitos} expresa esa elección en términos de requisitos medibles derivados
del propio sistema, en lugar de una preferencia general por la plataforma.

\begin{table}[ht!]
\caption{Requisitos de la plataforma de procesado, derivados del sistema implementado.}
\label{table:requisitos}
\centering
\begin{tabular}{|p{4.1cm}|p{2.4cm}|p{5.7cm}|}
\hline
\textbf{Requisito} & \textbf{Magnitud} & \textbf{Origen} \\
\hline
Entrada I2S nativa & --- & El micrófono es digital. Sin este periférico haría falta un
conversor externo, con la sensibilidad al ruido electromagnético que se pretendía evitar \\
\hline
Coma flotante en hardware & 5 transformadas de 1024 puntos en \mbox{30 ms} & Cada ráfaga
exige tres transformadas de vibración y dos de audio. La cifra corresponde a la medida en
la placa \\
\hline
Memoria de datos & \mbox{22,9 KiB} en reserva estática & Almacenes de ráfaga, de audio y de
trabajo de la transformada, dimensionados en tiempo de compilación \\
\hline
Conectividad inalámbrica integrada & \mbox{2,4 GHz} & Enlace con el concentrador. La banda
de \mbox{5 GHz} no es utilizable por el nodo \\
\hline
Capacidad de inferencia embarcada & --- & Requisito derivado del planteamiento del trabajo:
el diagnóstico debe poder ejecutarse sin el concentrador \\
\hline
\end{tabular}
\end{table}

Los dos primeros son los que efectivamente discriminan. La ausencia de entrada I2S nativa obliga
a incorporar un conversor analógico-digital externo, lo que reintroduce precisamente la
vulnerabilidad a la interferencia electromagnética que motivó la elección de un micrófono digital
en un entorno con un motor de inducción; la ausencia de unidad de coma flotante en hardware no
impide el cálculo pero obliga a reformular la transformada en aritmética de punto fijo, con el
consiguiente análisis de escalado y pérdida de precisión.

La reserva estática de \mbox{22,9 KiB} no es el consumo total del firmware sino solo el de los
almacenes de señal. Dimensionarlos en tiempo de compilación, en lugar de asignarlos de forma
dinámica, responde a que el nodo debe operar de forma continuada durante días: la asignación
dinámica repetida fragmenta el montón de memoria y termina provocando fallos que se manifiestan
tras horas de funcionamiento, cuando resultan más difíciles de diagnosticar.

Resuelta la plataforma, queda por fijar la cadencia de captura, que condiciona por completo el
análisis posterior y constituye la decisión técnica más restrictiva del sistema. El teorema del
muestreo establece
que para observar una componente de frecuencia $f$ es necesario muestrear por encima de
$2f$. Un compresor hermético accionado por un motor de inducción de dos polos alimentado a
\mbox{50 Hz} gira en torno a las \mbox{2900 RPM}, es decir, unas 48 vueltas por segundo, de
modo que su vibración presenta una componente fundamental próxima a los \mbox{48 Hz} junto
con sus armónicos.

Una cadencia de adquisición de \mbox{1 Hz}, adecuada para las magnitudes térmicas, sitúa el
límite de Nyquist en \mbox{0,5 Hz} y resulta por tanto dos órdenes de magnitud insuficiente
para la vibración. El efecto no es una medida degradada sino una medida carente de sentido
físico: la componente de \mbox{48 Hz} se repliega sobre la banda observable y aparece como
una frecuencia arbitraria sin relación con el régimen real del activo.

Transmitir de forma continua los tres ejes a la frecuencia requerida supondría del orden de
3000 valores por segundo, volumen incompatible con el enlace inalámbrico y contrario al
planteamiento del trabajo. La solución adoptada consiste en desacoplar las dos escalas
temporales del problema mediante dos canales de adquisición y publicación independientes,
según recoge el cuadro~\ref{table:muestreo}.

\begin{table}[ht!]
\caption{Canales de adquisición y sus parámetros de muestreo.}
\label{table:muestreo}
\centering
\begin{tabular}{|p{3.1cm}|p{2.2cm}|p{2.2cm}|p{4.6cm}|}
\hline
\textbf{Canal} & \textbf{Cadencia} & \textbf{Muestreo} & \textbf{Contenido} \\
\hline
\texttt{fridge/sensors} & \mbox{1 Hz} continuo & \mbox{1 Hz} & Nueve variables
instantáneas: temperaturas, aceleración, giro y nivel acústico \\
\hline
\texttt{fridge/vibration} & Una ráfaga cada \mbox{30 s} & \mbox{1 kHz} (vibración),
\mbox{16 kHz} (audio) & Características temporales y espectrales extraídas de 1024
muestras por eje \\
\hline
\end{tabular}
\end{table}

La ráfaga de 1024 muestras a \mbox{1 kHz} abarca \mbox{1,024 s} de señal, sitúa la
frecuencia de Nyquist en \mbox{500 Hz} y proporciona una resolución espectral de
\mbox{0,98 Hz}, suficiente para resolver la componente fundamental del compresor y varios
de sus armónicos. El filtro interno del acelerómetro se configura a \mbox{260 Hz} para
acotar el contenido por debajo de la frecuencia de Nyquist y evitar el replegado de las
componentes altas. Conviene precisar que este filtro, y no el criterio de Nyquist, es lo que
acota en la práctica la banda aprovechable: el contenido por encima de \mbox{260 Hz} llega
atenuado, de modo que la banda útil alcanza hasta el quinto armónico de la componente
fundamental.

El bus I\textsuperscript{2}C opera a \mbox{200 kHz}. La velocidad se fijó de forma
experimental: a \mbox{400 kHz} aparecían muestras corruptas de forma recurrente, atribuibles
a la integridad de la señal en un cableado sometido a vibración, y reducirla a la mitad
eliminó ese modo de fallo. A \mbox{200 kHz} la lectura de los seis bytes correspondientes a
los tres ejes consume unos \mbox{420 $\mu$s}, holgadamente dentro del presupuesto de
\mbox{1000 $\mu$s} por muestra.

Cabe señalar que el canal lento se mantuvo sin modificaciones al introducir el canal de
ráfaga, con objeto de conservar la resolución temporal que requieren las magnitudes térmicas
y la identificación de los ciclos de marcha y parada del activo. Los contadores de salud
descritos más adelante constituyen la única ampliación posterior de su formato.

\section{Extracción de características en el nodo}

Las características se calculan en el propio microcontrolador y solo ellas se transmiten, en
lugar de la señal en bruto. Sobre cada eje se obtienen el valor eficaz de la componente
alterna, el valor de pico, la kurtosis y la frecuencia y amplitud de la componente
dominante; sobre la señal acústica se obtiene el reparto de energía en cuatro bandas
espectrales. Esta reducción disminuye el volumen transmitido en más de tres órdenes de
magnitud y materializa el desplazamiento de la lógica de análisis al extremo de la red que
persigue el trabajo.

El cálculo espectral aplica una ventana de Hann antes de la transformada, con objeto de
reducir la fuga espectral de una señal que no es periódica en la ventana de análisis. La
estimación del pico se refina mediante interpolación parabólica sobre la logaritmo-magnitud
de los tres coeficientes centrales. El cuadro~\ref{table:estimador} recoge la mejora
obtenida, medida sobre señales sintéticas de frecuencia y amplitud conocidas en el intervalo
comprendido entre \mbox{45 Hz} y \mbox{97 Hz}.

\begin{table}[ht!]
\caption{Error del estimador espectral en el peor caso, sobre señales sintéticas.}
\label{table:estimador}
\centering
\begin{tabular}{|p{5.5cm}|p{3.2cm}|p{3.2cm}|}
\hline
\textbf{Estimador} & \textbf{Error en frecuencia} & \textbf{Error en amplitud} \\
\hline
Coeficiente de máxima magnitud & \mbox{0,49 Hz} & 13,2 \% \\
\hline
Interpolación parabólica & \mbox{0,016 Hz} & 3,6 \% \\
\hline
\end{tabular}
\end{table}

La interpolación resuelve la frecuencia unas sesenta veces por debajo del ancho del
coeficiente espectral, lo que permite seguir desplazamientos pequeños del régimen de giro,
y acota el error de amplitud debido a la pérdida de festoneado de la ventana. Las rutinas de
extracción se implementaron sin dependencias del entorno de ejecución del microcontrolador,
de modo que pueden verificarse numéricamente en un computador frente a señales de
propiedades conocidas analíticamente antes de su despliegue en la placa.

\section{Protocolo y modelo de intercambio de datos}

La comunicación entre el nodo y el concentrador emplea MQTT, estándar
ISO/IEC~20922~\cite{isoiec20922}, sobre TCP/IP y red inalámbrica
IEEE~802.11~\cite{ieee80211} en la banda de \mbox{2,4 GHz}. Se eligió MQTT por su sobrecarga reducida frente a alternativas
basadas en HTTP y por su modelo de publicación-suscripción, que desacopla el nodo de los
consumidores de datos.

Cada ciclo de medida genera un objeto JSON con nueve variables físicas, publicado en el tema
\texttt{fridge/sensors} con calidad de servicio~0. Se aceptó esta calidad de servicio
porque el análisis opera sobre ventanas estadísticas y la pérdida ocasional de una muestra
no altera el resultado, mientras que los mecanismos de confirmación introducirían
sobrecarga en el nodo.

El cuadro~\ref{table:variables} recoge las nueve variables del canal lento con su unidad y
su intervalo esperado. Los valores de los intervalos corresponden a la configuración del
firmware: fondo de escala de $\pm$\mbox{4 g} en el acelerómetro y de
$\pm$\mbox{500 $^\circ$/s} en el giroscopio.

\begin{table}[ht!]
\caption{Variables del canal lento, con su unidad e intervalo esperado.}
\label{table:variables}
\centering
\begin{tabular}{|p{2.1cm}|p{1.7cm}|p{2.4cm}|p{6.2cm}|}
\hline
\textbf{Variable} & \textbf{Unidad} & \textbf{Intervalo} & \textbf{Origen y observaciones} \\
\hline
\texttt{tempExt} & $^\circ$C & $-10$ a 40 & Sonda DS18B20. El valor $-127$ señala sonda
ausente o fallo de lectura \\
\hline
\texttt{accX/Y/Z} & m/s$^2$ & $\pm$39,2 & Acelerómetro MPU-6050. El eje alineado con la
vertical mide la gravedad en reposo \\
\hline
\texttt{gyroX/Y/Z} & rad/s & $\pm$8,7 & Giroscopio MPU-6050 \\
\hline
\texttt{motorTemp} & $^\circ$C & 15 a 85 & Temperatura interna del MPU-6050. Indicador
relativo de la superficie sobre la que se fija, no medida calibrada del motor \\
\hline
\texttt{noise} & --- & 0 a $10^5$ & Nivel acústico relativo del INMP441. No expresa
presión sonora en dB \\
\hline
\end{tabular}
\end{table}

La ausencia de dato se señala con el valor $-127$ en la temperatura externa y con un bloque de
ceros en las siete variables inerciales. La distinción importa porque un cero es medida válida en
casi cualquier variable, pero no en el eje que soporta la gravedad, donde \mbox{0,00 m/s$^2$}
resulta físicamente imposible con el sensor operativo y sirve por tanto de centinela sin
ambigüedad; el registrador escribe los campos ausentes como celdas vacías y nunca como ceros, por
el mismo motivo. El desglose completo de las tramas figura en el anexo~\ref{anexo:software}.

\section{Integridad del dato frente al fallo del bus}

El bucle principal no emplea esperas activas prolongadas: las dos cadencias se gobiernan
comparando el reloj de milisegundos del microcontrolador con el instante de la última publicación
de cada canal, y la conversión de la sonda de temperatura —unos \mbox{750 ms}— se solicita en un
ciclo y se recoge en el siguiente. La captura de la ráfaga es la única excepción deliberada y
bloquea el bucle unos \mbox{1,02 s}, porque el análisis frecuencial exige muestras equiespaciadas
y cualquier trabajo intercalado introduciría ruido espectral.

La vibración del activo provoca microcortes en el cableado del acelerómetro, que es precisamente
la magnitud que el sistema debe medir. Esta circunstancia condiciona el diseño de la adquisición
más que ninguna otra. El planteamiento inicial confiaba en que la biblioteca de acceso al bus
señalase el fracaso de una transacción devolviendo menos bytes de los solicitados, y se comprobó
que no lo hace: el controlador registra el error internamente mientras la función devuelve la
cuenta pedida. El código leía entonces de un almacén vacío, cuya lectura devuelve $-1$, valor que
al asignarse a una variable sin signo de ocho bits se convierte en \texttt{0xFF}, de modo que los
tres ejes resultaban $-1$ y esa muestra se publicaba como medida válida.

El efecto es desproporcionado respecto a la magnitud del fallo: una única muestra anómala entre
las 1024 de una ráfaga elevó la kurtosis del eje vertical desde 2,9 hasta 847, en coincidencia
exacta con la predicción analítica para un impulso aislado. Y la gravedad no reside en la pérdida
de una medida sino en que \textit{el artefacto imita el fenómeno que el sistema debe detectar}: la
kurtosis se seleccionó precisamente por su sensibilidad a los impulsos aislados, indicador
temprano clásico del deterioro de un rodamiento, de modo que un impulso de origen eléctrico y un
impacto mecánico incipiente resultan indistinguibles para ese estadístico. El efecto sobre la
estimación espectral es equivalente, dado que un impulso aislado presenta espectro plano, domina
la transformada y desplaza la frecuencia dominante a un valor arbitrario. La corrupción de una
sola muestra inutiliza por tanto las dos características principales del sistema, y un modelo
entrenado sobre datos así contaminados consideraría normal un intervalo de valores tan amplio que
perdería toda capacidad de discriminación.

La corrección se articula en tres mecanismos, cuyo detalle figura en el
anexo~\ref{anexo:software}. El primero \textbf{valida cada transacción} en cuatro puntos y admite
un reintento único por muestra, con lo que solo dos fallos consecutivos —que no responden a un
microcorte sino a un bus caído— descartan la ráfaga completa. El segundo \textbf{recupera el bus}
reiniciándolo y reconfigurando el sensor sin reiniciar el microcontrolador, porque un reinicio
interrumpiría la continuidad temporal de la señal, que es la información que el análisis explota;
lo acompaña una verificación de plausibilidad física del módulo de la aceleración, cuyo límite
demostrado es que no detecta la caída de un solo eje. El tercero no corrige nada sino que
\textbf{hace medible la calidad} de cada medida, y su ausencia es la que más daño causa: un dato
perdido que no se contabiliza resulta indistinguible de uno correcto al analizar el conjunto.
Ningún descarte se produce por tanto en silencio, y el nodo publica cuatro contadores de salud
—ráfagas descartadas, tramas inválidas, y reintentos de la ráfaga en curso y acumulados— que
convierten la calidad de cada medida en magnitud observable. El último, al ser monótono creciente
por construcción, delata además un reinicio del microcontrolador, circunstancia que de otro modo
no dejaría rastro en el conjunto de datos.

\section{Diseño del motor de análisis}

El cuadro~\ref{table:caracteristicas} recoge las características que el nodo extrae de cada
ráfaga, con la interpretación diagnóstica de cada una. La selección responde a la práctica
establecida en monitorización de condición de máquina rotativa y no a un criterio propio.

\begin{table}[ht!]
\caption{Características extraídas por eje y su interpretación diagnóstica.}
\label{table:caracteristicas}
\centering
\begin{tabular}{|p{2.3cm}|p{1.7cm}|p{8.4cm}|}
\hline
\textbf{Característica} & \textbf{Unidad} & \textbf{Interpretación} \\
\hline
Valor eficaz & m/s$^2$ & Nivel global de vibración. Indicador de severidad, de evolución
monótona con la degradación \\
\hline
Valor de pico & m/s$^2$ & Máximo absoluto tras eliminar la componente continua. Sensible a
impactos que el valor eficaz promedia \\
\hline
Kurtosis & --- & Forma de la distribución. El valor 3 corresponde a ruido gaussiano y 1,5 a
una senoide pura; los valores elevados indican impulsividad, indicador temprano del
deterioro de rodamientos \\
\hline
Frecuencia dominante & Hz & Régimen de giro. La aparición del segundo o tercer armónico como
componente dominante apunta a aflojamiento o desalineación \\
\hline
Amplitud dominante & m/s$^2$ & Amplitud en la frecuencia dominante. Distingue el aumento
general de vibración del aumento localizado en la frecuencia de giro, propio del
desequilibrio de masa \\
\hline
\end{tabular}
\end{table}

El par formado por el valor eficaz y la kurtosis es el que aporta capacidad de
discriminación. El primero cuantifica la magnitud y el segundo la naturaleza del fenómeno:
un desequilibrio de masa eleva el valor eficaz manteniendo una kurtosis próxima a la de una
senoide, mientras que un rodamiento en deterioro incipiente apenas altera el valor eficaz y
sin embargo dispara la kurtosis. Son mecanismos de fallo distintos que ese par separa.

Procede exponer a continuación la estrategia de detección. La detección se plantea como un problema \textit{no supervisado} de una sola clase. La razón
es la disponibilidad de datos: de un activo en servicio se obtiene con facilidad su
comportamiento normal, mientras que los ejemplos de fallo son escasos y deben provocarse de
forma deliberada. Un clasificador supervisado exigiría un número de ejemplos de cada modo de
fallo que no es alcanzable en este contexto.

El procedimiento previsto consta por tanto de dos etapas asimétricas. El modelo se ajusta
sobre el comportamiento nominal registrado y aprende a describir esa región del espacio de
características; en explotación, toda observación suficientemente alejada de ella se marca
como anómala. El umbral se fija sobre un percentil de la distribución de puntuaciones del
propio conjunto nominal, lo que hace explícita y ajustable la tasa de falsos positivos
admitida.

Dos decisiones de diseño acompañan a este planteamiento. La primera es que el entrenamiento
se realiza \textit{por lotes} y fuera del nodo, no de forma incremental. Un modelo que se
reajustase de forma continua sobre las observaciones recientes acabaría incorporando la
degradación progresiva del activo a su noción de normalidad, con lo que dejaría de señalarla:
el mecanismo de detección se anularía a sí mismo justamente en el escenario que debe
detectar. El reajuste queda por ello condicionado a una intervención humana que confirme el
estado del activo, criterio análogo al de la recalibración de instrumentación tras una
reparación.

La segunda es que en el nodo no se embarca el modelo completo sino un estadístico equivalente
derivado de él, del orden de unos pocos kilobytes: el vector de medias y la dispersión de
cada característica sobre el estado nominal, junto con el umbral. Esta reducción permite
ejecutar la inferencia sin una infraestructura de aprendizaje automático en el
microcontrolador, y ofrece además una vía de verificación: pasar un mismo conjunto de
ráfagas por el modelo completo y por el estadístico embarcado permite cuantificar la
concordancia entre ambos y, con ella, el coste de la simplificación.

El algoritmo concreto se seleccionó comparando siete candidatos sobre los datos disponibles,
con arreglo a los criterios que se enuncian a continuación y cuyo resultado se recoge en el
capítulo~\ref{cap:resultados}. Procede justificar el orden en que se aplican, porque no es el
habitual.

El criterio principal \textbf{no} es la tasa de detección. Con un conjunto nominal reducido a
un episodio de marcha, la tasa de detección se satura para cualquier candidato razonable y por
tanto no distingue entre ellos. Lo que sí distingue es la \textbf{estabilidad de la frontera de
decisión}: se ajusta cada modelo sobre submuestras de tamaño creciente del conjunto nominal y
se mide la dispersión de la tasa de falsos positivos fuera de muestra. Un modelo cuya frontera
se desplaza al cambiar qué observaciones la ajustan no es utilizable con los datos disponibles,
por elevada que sea su tasa de acierto.

El segundo criterio es la tasa de falsos positivos \textbf{fuera de muestra}, contrastada con
el valor que el umbral fija por construcción. La tasa dentro de muestra la determina el propio
umbral y no aporta información.

El cuarto criterio, tras la tasa de detección, es el \textbf{número de hiperparámetros}. Cada
grado de libertad que el conjunto de datos no permite ajustar honestamente no constituye
capacidad del modelo sino riesgo de sobreajuste, y con observaciones correlacionadas
procedentes de un episodio no hay margen para ajustar ninguno.

El quinto es la \textbf{viabilidad de embarcar la inferencia}, que en este trabajo no es una
consideración de implementación sino el objetivo planteado: un modelo que exija el
concentrador en funcionamiento no resuelve el problema enunciado en el
capítulo~\ref{cap:objetivos}.

De esta ordenación se sigue una consecuencia que conviene anticipar: entre candidatos con
comportamiento equivalente, la elección recae en el más simple. No por economía de medios, sino
porque con el tamaño de muestra disponible la simplicidad es la propiedad que hace el resultado
verificable.

El procedimiento completo, con los umbrales declarados de forma explícita, está en
\texttt{server/analisis/}, cuyo fichero \texttt{README.md} recoge la justificación de cada
decisión de limpieza y de derivación de características.

\section{Los modelos considerados y el criterio de elección}

Establecido que el problema es de detección de una clase, procede enunciar qué familias de
modelos resultan aplicables y en qué se diferencian, dado que la elección entre ellas condiciona
tanto la capacidad de detección como el coste de embarcar la inferencia.

Se consideraron cinco modelos y dos reglas de decisión. Los cuatro primeros se designan por su
nombre habitual en la literatura, que es el inglés; se acompaña la traducción al castellano en
su primera aparición y en lo sucesivo se emplea la denominación original, todos ellos ajustables únicamente sobre
observaciones del estado nominal. El cuadro~\ref{table:modelos} resume su fundamento.

\begin{table}[ht!]
\caption{Modelos de detección de una clase considerados.}
\label{table:modelos}
\centering
\begin{tabular}{|p{3.6cm}|p{10.4cm}|}
\hline
\textbf{Modelo} & \textbf{Fundamento} \\
\hline
\textit{Elliptic Envelope}\newline\small(envolvente elipsoidal robusta) & \normalsize Ajusta un elipsoide a la nube de observaciones nominales y mide la
distancia de Mahalanobis al centro. La estimación de la covarianza es robusta, a partir del
subconjunto de menor determinante, para que unas pocas observaciones atípicas del propio
conjunto de ajuste no desplacen la frontera~\cite{rousseeuw1999} \\
\hline
\textit{Local Outlier Factor}\newline\small(factor local de atipicidad) & \normalsize Compara la densidad local de cada observación con la de su
vecindad y declara atípico lo que reside en una región menos poblada que su entorno inmediato.
No presupone forma alguna para la distribución nominal~\cite{breunig2000} \\
\hline
\textit{Isolation Forest}\newline\small(bosque de aislamiento) & \normalsize Construye árboles con particiones aleatorias y mide cuántas hacen falta
para aislar cada observación. Las atípicas se aíslan antes, con un coste que no depende de la
dimensión~\cite{liu2008} \\
\hline
\textit{One-Class SVM}\newline\small(vectores soporte de una clase) & \normalsize Determina la frontera de menor volumen que encierra la mayor
parte de las observaciones nominales, en un espacio transformado por un núcleo~\cite{scholkopf2001} \\
\hline
Envolvente de Mahalanobis & La misma forma cuadrática, con estimación clásica de la covarianza
y regularización diagonal. Se incluye como referencia de mínimo coste \\
\hline
Regla sobre el número de picos & Umbral sobre una única característica: cuántos picos
espectrales superan una fracción de la amplitud mayor \\
\hline
Regla sobre la relación armónica & Intervalo sobre el cociente entre un pico secundario y la
fundamental del giro \\
\hline
\end{tabular}
\end{table}

Las dos últimas no son modelos de aprendizaje sino umbrales construidos a mano, y se incluyen
precisamente por eso: si una regla de dos comparaciones iguala a un modelo ajustado, la
complejidad adicional de este no está justificada.

La elección entre ellos no se resuelve por la tasa de detección, y conviene enunciar por qué antes de
exponer el procedimiento. Con un conjunto de referencia de tamaño reducido y un único fallo
disponible, la tasa de detección se satura para cualquier candidato razonable y por tanto no
discrimina entre ellos. Los criterios que sí discriminan son cuatro, aplicados en este orden.

El primero es el \textbf{comportamiento sobre condiciones de operación no observadas}, estimado
dejando fuera un episodio de marcha completo y midiendo los falsos positivos sobre él. Es el
criterio decisivo, y exige que la campaña de referencia comprenda un número suficiente de
episodios.

El segundo es la \textbf{tasa de falsos positivos} contrastada con el valor que el umbral fija
por construcción, dado que la tasa medida sobre el propio conjunto de ajuste la determina ese
umbral y no informa de nada.

El tercero es el \textbf{número de hiperparámetros}. Cada grado de libertad que el conjunto de
datos no permite ajustar honestamente no constituye capacidad del modelo sino riesgo de
sobreajuste.

El cuarto es el \textbf{coste de embarcar la inferencia}, que en este trabajo no es una
consideración de implementación sino el objetivo planteado en el capítulo~\ref{cap:objetivos}.
Debe consignarse, no obstante, que este criterio se empleó inicialmente para descartar
candidatos \textit{sin haberlo medido}, y que la medición posterior recogida en el
capítulo~\ref{cap:resultados} demostró que la restricción no existía.

De esta ordenación se sigue que entre candidatos equivalentes la elección recae en el más
simple, no por economía de medios sino porque con este tamaño de muestra la simplicidad es la
propiedad que hace el resultado verificable.

\section{Componente tonal de alta frecuencia y acotación de la banda útil}

Al reubicar el acelerómetro en un punto de transmisión mecánica directa sobre la carcasa del
compresor (anexo~\ref{anexo:hardware}), el nivel de vibración medido se multiplicó por un factor
entre 9 y 20, dejando al descubierto una componente tonal dominante en torno a \mbox{448 Hz}. Se
atribuyó en un primer momento a una \textbf{resonancia mecánica del montaje adhesivo}, por dos
indicios: aparecía ya con el sensor en la posición anterior, aunque con amplitud del orden del
ruido, y se desplazaba entre \mbox{398} y \mbox{448 Hz} durante la sesión, comportamiento
incompatible con un artefacto eléctrico, que se habría mantenido fijo. El análisis expuesto en el
capítulo~\ref{cap:resultados} refutó esa atribución: el desplazamiento aparente no era la deriva de
una resonancia sino \textbf{dos armónicos distintos} observados en ráfagas distintas, y lo que hay
es una tríada de armónicos del giro del motor —\mbox{398}, \mbox{448} y \mbox{497 Hz}, los órdenes
8, 9 y 10 con desviación inferior al 0,05~\% respecto del entero.

Esta componente tonal dominaba el espectro y fijaba la kurtosis temporal en torno a 1,72 (próxima a
los 1,5 de una senoide pura), enmascarando el régimen de giro fundamental del motor (\mbox{49,77 Hz})
y limitando la sensibilidad ante impactos transitorios. Dado que el nodo procesa y descarta la señal
bruta en el extremo, la discriminación debía resolverse en el propio firmware antes de transmitir:

\begin{enumerate}
    \item \textbf{Picos espectrales múltiples ordenados:} En lugar del pico dominante único, el nodo
        extrae y publica los tres picos de mayor amplitud por eje con una guarda de cuatro
        coeficientes de la transformada, capturando simultáneamente la fundamental y los armónicos dominantes.
    \item \textbf{Filtrado paso bajo a \mbox{150 Hz}:} Se aplica un filtro digital paso bajo de
        segundo orden antes de calcular los estadísticos temporales —valor eficaz, pico y kurtosis—, aislando
        la banda de la fundamental y dejando el cálculo espectral sin filtrar hasta los \mbox{500 Hz}
de Nyquist. El corte se derivó de la práctica de no emplear un montaje por encima de un
        tercio de su frecuencia de resonancia, aplicada a los \mbox{448 Hz} que entonces se
        atribuían al acoplamiento; refutada esa atribución, queda como cota conservadora sin
        verificación experimental, según se declara en el capítulo~\ref{cap:resultados}.
\end{enumerate}

El firmware publica los tres picos \textbf{ordenados por amplitud} y no por frecuencia, de modo
que la derivación de características ha de reordenarlos: en caso contrario un mismo fenómeno
arroja cocientes incoherentes cuando un armónico supera a la fundamental y le arrebata la
posición de dominante. El capítulo~\ref{cap:resultados} recoge el defecto que esa
circunstancia introdujo.

El cuadro~\ref{table:resonancia-efecto} resume el efecto medido de ambas medidas sobre la misma señal.

\begin{table}[ht!]
\caption{Efecto de las correcciones sobre las características del eje vertical.}
\label{table:resonancia-efecto}
\centering
\begin{tabular}{|p{4.2cm}|p{3.4cm}|p{4.4cm}|}
\hline
\textbf{Característica} & \textbf{Antes} & \textbf{Después} \\
\hline
Valor eficaz & \mbox{0,35 m/s$^2$} & \mbox{0,062 m/s$^2$}, banda \mbox{0-150 Hz} \\
\hline
Kurtosis & 1,72 (fija) & 2,48 a 7,35 \\
\hline
Fundamental del activo & No identificable & \mbox{49,77 Hz}, tercer pico, en 6 de 7 ráfagas \\
\hline
\end{tabular}
\end{table}

\section{Estándares aplicados y consideraciones de seguridad}

El sistema se apoya en MQTT (ISO/IEC~20922)~\cite{isoiec20922} e IEEE~802.11~\cite{ieee80211}
para la comunicación, con seguridad WPA/WPA2 en el enlace inalámbrico. El intermediario de
mensajería opera sin cifrado de transporte ni autenticación de cliente, configuración admisible
para una prueba de concepto en una red de laboratorio no expuesta a Internet.

Durante el desarrollo del firmware, y de forma transitoria, el nodo publicó contra un
intermediario público accesible desde Internet, empleando un prefijo de tema aleatorio. Procede
señalar que dicho prefijo \textbf{no constituía una medida de seguridad}: cualquier tercero que
lo conociera podría publicar tramas en los mismos temas, que el registrador incorporaría al
conjunto de datos como si procedieran del nodo. La consecuencia no sería una pérdida de
confidencialidad sino un riesgo de integridad capaz de invalidar el ajuste del modelo sin
manifestarse. Se estableció por ello como criterio que ninguna campaña destinada a la memoria se
realizase en esa configuración, abandonada al desplegar el concentrador.

En la configuración definitiva el intermediario reside en el concentrador, que genera además la
red inalámbrica sin conexión con ninguna otra. El aislamiento desplaza la garantía de integridad
a la clave de acceso de esa red, circunstancia cuya implicación resulta poco intuitiva: en este
sistema la robustez de una contraseña no protege la privacidad de un dato sino \textbf{la validez
de un resultado experimental}.

Un despliegue en explotación requeriría autenticación de cliente y cifrado TLS del canal. Su
ausencia se justifica por el alcance de prueba de concepto y por el aislamiento de la red, y no
por el coste o la capacidad del nodo, que soporta TLS.
